\chapter{Establishing API Communication}

This chapter connects everything from Chapters 1--7 into one fully working
round trip: a React page calling \texttt{api.get("/tasks")} and rendering
whatever MongoDB returns. Read this chapter slowly -- every CRUD feature for
the rest of the book is a variation of this exact pattern.

\section{The Complete Cycle, Labeled}

\begin{diagrambox}[GET /api/tasks -- Full Cycle]
\begin{verbatim}
 React: Dashboard.jsx
   useEffect(() => { taskService.getTasks() }, [])
        |
        v
 services/taskService.js
   api.get("/tasks")
        |
        v
 Axios instance (src/api/axios.js)
   + baseURL:  http://localhost:5000/api
   + withCredentials: true (sends JWT cookie)
        |
        v
 HTTP Request  ------------------------------>
   GET /api/tasks
   Headers: Cookie: token=<jwt>; Content-Type: application/json
   Body: (none - GET has no body)
        |
        v
 Express app.js
   app.use("/api/tasks", taskRoutes)
        |
        v
 routes/taskRoutes.js
   router.get("/", protect, getTasks)
        |
        v
 middleware/auth.js  (protect)
   verifies JWT from cookie -> attaches req.user
        |
        v
 controllers/taskController.js
   getTasks(req, res)
        |
        v
 Mongoose
   Task.find({ user: req.user._id })
        |
        v
 MongoDB
   scans "tasks" collection, returns matching documents
        |
        v
 controllers/taskController.js
   res.status(200).json({ success: true, data: tasks })
        |
        v
 HTTP Response  <------------------------------
   Status: 200 OK
   Body: { "success": true, "data": [ {...}, {...} ] }
        |
        v
 Axios instance
   resolves promise with response.data
        |
        v
 services/taskService.js
   returns res.data.data  (the array of tasks)
        |
        v
 React: Dashboard.jsx
   setTasks(data) -> re-render -> TaskList shows TaskCards
\end{verbatim}
\end{diagrambox}

\section{The Frontend Side}

\begin{lstlisting}[language=JavaScript]
// src/services/taskService.js
import api from "../api/axios";

export async function getTasks() {
  try {
    const res = await api.get("/tasks");
    return res.data.data;
  } catch (err) {
    throw new Error(err.response?.data?.message || "Failed to load tasks");
  }
}
\end{lstlisting}

\begin{lstlisting}[language=JavaScript]
// src/pages/Dashboard.jsx
import { useEffect, useState } from "react";
import { getTasks } from "../services/taskService";
import TaskList from "../components/TaskList";

export default function Dashboard() {
  const [tasks, setTasks] = useState([]);
  const [error, setError] = useState("");

  useEffect(() => {
    getTasks()
      .then(setTasks)
      .catch((err) => setError(err.message));
  }, []);

  if (error) return <p className="error">{error}</p>;
  return <TaskList tasks={tasks} />;
}
\end{lstlisting}

\section{The Backend Side}

\begin{lstlisting}[language=JavaScript]
// routes/taskRoutes.js
const express = require("express");
const router = express.Router();
const { protect } = require("../middleware/auth");
const { getTasks, createTask } = require("../controllers/taskController");

router.get("/", protect, getTasks);
router.post("/", protect, createTask);

module.exports = router;
\end{lstlisting}

\begin{lstlisting}[language=JavaScript]
// controllers/taskController.js
const Task = require("../models/Task");

exports.getTasks = async (req, res) => {
  try {
    const tasks = await Task.find({ user: req.user._id }).sort({ createdAt: -1 });

    res.status(200).json({
      success: true,
      data: tasks,
    });
  } catch (err) {
    res.status(500).json({
      success: false,
      message: "Failed to fetch tasks",
    });
  }
};
\end{lstlisting}

\begin{lstlisting}[language=JavaScript]
// app.js (wiring)
const express = require("express");
const app = express();

app.use(express.json());
app.use("/api/tasks", require("./routes/taskRoutes"));
app.use("/api/auth", require("./routes/authRoutes"));

module.exports = app;
\end{lstlisting}

\begin{notebox}[NOTE]
Notice the controller never talks to Express's \texttt{req}/\texttt{res} objects
inside Mongoose calls, and never talks to MongoDB directly inside routing
code. Each layer only knows about the layer immediately next to it -- that's
what makes this chain easy to debug and to test.
\end{notebox}

\section{Standard Response Shapes}

To keep the frontend's unwrapping logic (\texttt{res.data.data}) consistent
everywhere, every endpoint in this book returns one of the following shapes.

\begin{table}[h]
\centering
\begin{tabularx}{\textwidth}{lXl}
\toprule
\textbf{Shape} & \textbf{Used for} & \textbf{Example} \\
\midrule
\texttt{\{ success, data \}} &
Single resource or list fetch/create/update &
\texttt{GET /tasks}, \texttt{POST /tasks} \\[4pt]
\texttt{\{ success, message \}} &
Actions with no payload to return &
\texttt{DELETE /tasks/:id}, \texttt{POST /auth/logout} \\[4pt]
\texttt{\{ success, data, page, limit, total, totalPages \}} &
Paginated list endpoints (Chapter 20) &
\texttt{GET /tasks?page=2\&limit=10} \\
\bottomrule
\end{tabularx}
\end{table}

\begin{warnbox}[WARNING]
Never mix shapes for the same endpoint across success and error cases --
e.g. returning \texttt{data} on success but a bare string on failure. Always
keep \texttt{success: boolean} as the first field so the frontend can branch on
it consistently regardless of HTTP status code edge cases.
\end{warnbox}

\begin{tipbox}[TIP]
When a new feature isn't working, check these five things in order: (1) is
the URL in \texttt{taskService.js} correct, (2) is the HTTP method correct, (3)
does the Express route match that method and path, (4) does the controller
query the right field, (5) does the response shape match what the frontend
expects. This covers the overwhelming majority of ``it's just not working''
bugs in a MERN app.
\end{tipbox}
